\section{研究问题与应用版图：Stone Lab 的统一问题}

Stone 给出的长期研究问题是：
\textit{``To what degree can autonomous intelligent agents learn in the presence of teammates and/or adversaries in real-time dynamic domains?''}

这个问题从 1998 年延续至今，说明 ``Agent'' 在学术上并不是新词，而是不断被新计算范式激活的老问题。

\begin{importantbox}{统一问题的四个关键词}
\begin{itemize}
    \item \textbf{learn}：不是手工脚本，而是策略学习；
    \item \textbf{teammates/adversaries}：多体互动不是附加项，而是主问题；
    \item \textbf{real-time}：控制频率和延迟约束是硬条件；
    \item \textbf{dynamic domains}：状态分布持续变化，离线最优策略会失效。
\end{itemize}
\end{importantbox}

围绕这一问题，Stone 的例子覆盖 robotics、robot soccer、autonomous driving、Gran Turismo、human feedback learning、ad hoc teamwork、ethical dataset construction。看上去分散，但都在回答同一个东西：\textbf{能否在复杂互动环境中持续学习并保持可用行为。}

\begin{warningbox}{研究叙事中的陷阱}
如果只看单一 benchmark 提升，容易误判为能力突破；但在动态多体环境中，很多 ``提分'' 只是对静态分布的过拟合。
\end{warningbox}

\subsection{本章小结}
统一问题让我们避免被技术热点牵着走。它同时约束了算法目标、系统设计和评测方式。

